\section{Concluzii}
În această lucrare am proiectat și implementat o funcție hash bazată pe sisteme dinamice discrete haotice, construită în jurul structurii Sponge și al arborelui Merkle.
Analiza statistică prezentată în capitolul anterior confirmă faptul că funcția propusă posedă proprietăți criptografice solide.
Testele de confuzie și difuzie arată că o modificare minimă a mesajului de intrare determină o schimbare semnificativă a hash-ului.
Rezistența la coliziuni este, de asemenea, foarte bună: distribuția numărului de bytes egali și suma diferențelor absolute ale bytes-ilor se apropie de valorile teoretice, iar numărul de coliziuni observate este sub limita teoretică.
În plus, valorile hash sunt distribuite uniform pe întreg spațiul de ieșire, ceea ce împiedică atacurile bazate pe criptanaliză statistică.
Toate aceste rezultate sunt comparabile cu cele raportate de funcțiile hash consacrate din literatura de specialitate, ceea ce indică un nivel de securitate adecvat aplicațiilor criptografice moderne. \\

Pe lângă proprietățile criptografice, funcția hash propusă se remarcă prin viteza și eficiența sa pe sistemele cu mai multe nuclee.
Datorită arborelui Merkle, calculul hash-ului poate fi paralelizat eficient, viteza de procesare crescând aproape liniar cu numărul de nuclee alocate.
Pe sistemele moderne testate, funcția atinge viteze de procesare de ordinul sutelor de MB/s, comparabile cu cele ale funcțiilor hash utilizate la nivel de industrie, precum SHA-256 sau MD5.
Acest lucru este cu atât mai semnificativ cu cât funcțiile hash consacrate rulează pe un singur nucleu și beneficiază adesea de optimizări hardware dedicate. \\

Ca direcții viitoare de dezvoltare, o primă posibilitate ar fi implementarea funcției în limbajul C, în vederea obținerii unor viteze de procesare și mai mari.
O a doua direcție o constituie explorarea metodelor de optimizare hardware, prin care funcția ar putea fi accelerată suplimentar pe arhitecturi specializate.
În cele din urmă, o direcție promițătoare ar fi extinderea funcției hash într-o schemă completă de criptare și decriptare, valorificând proprietățile sistemelor dinamice haotice care stau la baza ei.